\Page{17}{\textenglish{04 / FUTURE}}{כשהחנות חוזרת לחדר הפיתוח}{38}{44}

{\fontsize{12.3}{17.9}\selectfont

חנות יכולה להמליץ לכל קונה על פריטים אחרים בלי לשנות אפילו הזמנה אחת מן המפעל. ההשפעה על הייצור מתחילה כשהצוות מבין מה אפשר ללמוד מן הבחירות האלה. לחיצה על גוון יכולה לתת כיוון. אם הרבה קונות מתעניינות במידה שאזלה, המכירות לבדן ממעיטות בהיקף הביקוש. ואם פריט שנמכר מהר חוזר בגלל התאמה, השיפור הדרוש עשוי להיות בגזרה ולא בכמות. הקשר בין התאמה אישית לתכנון ייצור תלוי בחיבור נתונים שלא נאספו במקור כדי לענות על אותה שאלה.

מחקר היסטורי על \textenglish{Zara} תיעד שימוש בתחזיות, במלאי לפי מידה ובכללי תצוגה כדי לבחור הקצאה בין חנויות. הוא מראה כיצד אפשר להפיק החלטת משלוח מתחזית תחת מגבלות, באמצעות אופטימיזציה. זאת דוגמה להחלטה על מלאי קיים. ההשפעה הסביבתית דורשת מדידה אחרת. הסוכנות האירופית לסביבה העריכה ב־2024 ש־4\% עד 9\% ממוצרי הטקסטיל שהוצעו בשוק האירופי מושמדים לפני ששימשו למטרתם. הטווח מבוסס על נתונים חלקיים, בלי לייחס את ההשמדה לתוכנה או לשיטת הזמנה מסוימת. \SourceRef{12}

\par}

\par\begingroup\color{Ink}\fontsize{11.5}{15.7}\selectfont
\renewcommand{\arraystretch}{1.45}
\setlength{\tabcolsep}{6pt}
\arrayrulecolor{Silver}\rowcolors{2}{Ivory}{PaperSteel}
\noindent\begin{tabularx}{\linewidth}{@{}>{\raggedleft\arraybackslash}p{.133\linewidth}*{3}{>{\raggedleft\arraybackslash}X}@{}}
\rowcolor{Graphite}
\textcolor{Ivory}{\textbf{החלטה}} & \textcolor{Ivory}{\textbf{תכנון מקובל}} & \textcolor{Ivory}{\textbf{תכנון בסיוע מודל}} & \textcolor{Ivory}{\textbf{מדד בדיקה}} \\
מה להזמין & שימוש בעונות קודמות ובמידע מהחנויות & תחזית לפי פריט, אזור ומידה & חוסרים לצד יתרת מלאי \\
מה לשנות & מעבר על החזרות והערות צוות & קיבוץ סיבות החזרה והפניה לבדיקת דוגמה & אותה סיבת החזרה אחרי התיקון \\
מה לעשות בעודף & מיון והחלטת מנהלת מלאי & הצעת העברה, תיקון או שינוי מלאי & כמה פריטים חזרו לשימוש ומה הושמד \\
\hline
\end{tabularx}\par
\endgroup


\vspace{3mm}\par\noindent
\begin{english}\centering
\begin{tikzpicture}[
  stage/.style={draw=Silver,fill=PaperSteel,align=center,text width=26mm,
    inner sep=1.25mm,minimum height=25mm,font=\fontsize{10.5}{13.5}\selectfont},
  flow/.style={draw=Warm,line width=.7pt,-{Latex[length=2mm]}}]
  \node[stage,fill=Graphite,draw=Graphite,text=Ivory] (data) at (14.4,0) {\texthebrew{ביקוש, מלאי והחזרות}};
  \node[stage] (forecast) at (10.8,0) {\texthebrew{תחזית וזיהוי חריגים}};
  \node[stage,fill=PaperPink] (decision) at (7.2,0) {\texthebrew{החלטת צוות על כמות, התאמה או חלוקה}};
  \node[stage] (operation) at (3.6,0) {\texthebrew{ייצור, הפצה ומכירה}};
  \node[stage,fill=Graphite,draw=Graphite,text=Ivory] (feedback) at (0,0) {\texthebrew{חוסרים, החזרות וגורל העודף}};
  \draw[flow] (data.west) -- (forecast.east);
  \draw[flow] (forecast.west) -- (decision.east);
  \draw[flow] (decision.west) -- (operation.east);
  \draw[flow] (operation.west) -- (feedback.east);
  \draw[flow,dashed] (feedback.south) -- (0,-1.8) -- (14.4,-1.8) -- (data.south);
\end{tikzpicture}\par
\end{english}






\Caption{הסכמה מימין לשמאל. נתוני ביקוש, מלאי והחזרות \textenglish{\(\leftarrow\)} תחזית וזיהוי חריגים \textenglish{\(\leftarrow\)} החלטת צוות על כמות, התאמה או חלוקה \textenglish{\(\leftarrow\)} ייצור, הפצה ומכירה \textenglish{\(\leftarrow\)} מדידת חוסרים, החזרות וגורל העודף. חץ משוב מחזיר תוצאות לנתונים.}

{\fontsize{12.3}{17.9}\selectfont

צמצום בזבוז מתחיל כשהסוף נראה לעין. אם אחרי שינוי בהזמנות נשארו פחות פריטים לא מכורים, צריך עדיין לברר מה קרה לחוסרים ומה היה גורלם של העודפים. רק אז אפשר לדעת אם ההחלטה הפחיתה בזבוז, או העבירה את הבעיה למקום אחר.

\par}
